\section{Explicit ordinal coordinates for the canonical invariant}

\begin{definition}[Chronological binary coordinate]
\label{hir:disjoint-invariant-ordinal-colex}
\lean{TranscendentalNakamaye.Pending.HironakaDisjointInvariantOrdinal.chronologicalCode,TranscendentalNakamaye.Pending.HironakaDisjointInvariantOrdinal.chronologicalCode_eq_sum}
\leanok
For a finite set $K$ of chronological divisor positions, put
$c(K)=\sum_{i\in K}2^i$. The actual positions themselves are used as exponents.
\end{definition}

\begin{lemma}[Bounded ordinal blocks]
\label{hir:disjoint-invariant-ordinal-blocks}
Suppose $r:E\longrightarrow\mathrm{Ord}$ is strictly increasing and
$r(x)<B$ for every $x$, with $0<B$. Then
$(n,x)\longmapsto Bn+r(x)$ is strictly increasing for the lexicographic
order on $\mathbb N\times E$, and every value is below $B\omega$.
\end{lemma}
\begin{proof}
For each $n$ and $x$, strict monotonicity of addition in its right argument
and distributivity give
$Bn+r(x)<Bn+B=B(n+1)<B\omega$; the last strict inequality follows from
$n+1<\omega$ and $B>0$. If $n<m$, then $n+1\le m$, so
$Bn+r(x)<B(n+1)\le Bm\le Bm+r(y)$. If $n=m$ and $x<y$, then
$Bn+r(x)<Bn+r(y)$, again by strict monotonicity in the right argument.
These are exactly the two cases of the lexicographic strict order.
\end{proof}

\begin{lemma}[Adding endpoints and summing blocks]
\label{hir:disjoint-invariant-ordinal-endpoints}
Suppose $r:E\longrightarrow\mathrm{Ord}$ and $s:F\longrightarrow\mathrm{Ord}$
are strictly increasing with positive bounds $B,C$. Adjoining a greatest
point to $E$, retaining $r$ and assigning $B$ to the new point, gives a
strictly increasing map bounded by $B+1$. If $B$ is a limit ordinal,
adjoining a least point, assigning it zero and assigning $r(x)+1$ to $x$,
gives a strictly increasing map bounded by $B$. The lexicographic sum
$E\mathbin{\oplus}F$ has strictly increasing coordinates $r(x)$ on the
left and $B+s(y)$ on the right, all below $B+C$.
\end{lemma}
\begin{proof}
Every old coordinate is strictly below the new greatest coordinate $B$,
and both kinds of coordinate lie below $B+1$. The new least coordinate
zero is strictly below each successor $r(x)+1$. Ordinal successor is
strictly increasing; because $B$ is a limit ordinal and $r(x)<B$, each
successor remains strictly below $B$. For the sum, a left coordinate is
below $B$, while each right coordinate is at least $B$. Comparisons
within the left component are unchanged; comparisons within the right
component are preserved by strict monotonicity of addition in its right
argument. Left coordinates lie below $B\le B+C$, and right coordinates
lie below $B+C$ by their bound $s(y)<C$.
\end{proof}

\begin{definition}[Recursive ordinal coordinates]
\label{hir:disjoint-invariant-ordinal-code}
\uses{hir:disjoint-invariant-ordinal-colex,hir:canonical-centres}
\lean{TranscendentalNakamaye.Pending.HironakaDisjointInvariantOrdinal.monomialCode,TranscendentalNakamaye.Pending.HironakaDisjointInvariantOrdinal.valueCode,TranscendentalNakamaye.Pending.HironakaDisjointInvariantOrdinal.valueBound}
\leanok
Write $\beta=(\omega\omega)\omega=\omega^3$ and put
$M(a,b,K)=(\omega\omega)a+\omega b+c(K)$.
Set $A_0=\omega$, $L_q=(A_q+1)\omega$,
$D_q=\beta+(L_q+1)$, and $A_{q+1}=D_q\omega$.
The original value type in dimension zero is $\mathbb N$, with coordinate
$r_0(n)=n$. To define the next coordinate on the original recursive type
\[
\mathbb N\times_{\mathrm{lex}}
\bigl((\{\bot\}\cup\operatorname{Val})\mathbin{\oplus}
(\mathbb N\times_{\mathrm{lex}}(\operatorname{ValTy}_q\cup\{\top\})
\cup\{\top\})\bigr),
\]
first put $t_q(v)=r_q(v)$ on finite lower values and $t_q(\top)=A_q$.
Set $u_q(s,v)=(A_q+1)s+t_q(v)$, and set $u_q(\top)=L_q$ for the
additional greatest point. On the left component set
$z_q(\bot)=0$ and $z_q(a,b,K)=M(a,b,K)+1$; on the right component set
$z_q(w)=\beta+u_q(w)$. Finally put
$r_{q+1}(o,w)=D_qo+z_q(w)$.
All sums and products here are ordinal operations in the displayed order.
\end{definition}

\begin{theorem}[Embedding of the original value type]
\label{hir:disjoint-invariant-ordinal-embedding}
\uses{hir:disjoint-invariant-ordinal-code,hir:disjoint-invariant-ordinal-blocks,hir:disjoint-invariant-ordinal-endpoints}
\lean{TranscendentalNakamaye.Pending.HironakaDisjointInvariantOrdinal.valueEmbedding,TranscendentalNakamaye.Pending.HironakaDisjointInvariantOrdinal.valueCode_lt_bound,TranscendentalNakamaye.Pending.HironakaDisjointInvariantOrdinal.valueCode_lt_iff}
\leanok
For each $q$, the recursive coordinate $r_q$ is an order embedding of the
original $\operatorname{ValTy}_q$ into the ordinals, and $r_q(v)<A_q$
for every value $v$.
\end{theorem}
\begin{proof}
Binary expansion is the inverse of the natural-number/colexicographic
finite-subset order isomorphism. Thus $c(K)$ preserves strict colex order,
and its ordinal cast is below $\omega$. Apply the bounded-block lemma
first to $(b,K)$ with bound $\omega$, then to $(a,(b,K))$ with bound
$\omega\omega$. The result is exactly $M$, strictly increasing and
bounded by $\beta$. The ordinal $\beta$ is a limit ordinal because it is
a positive ordinal multiplied by $\omega$.
The base map $n\mapsto n$ is strictly increasing and bounded by $\omega$.
Assume inductively that $r_q$ is strictly increasing and bounded by the
positive ordinal $A_q$. The greatest-point construction gives $t_q$ with
bound $A_q+1$. The block lemma gives $u_q$ with bound $L_q$, and another
greatest-point construction extends its domain with bound $L_q+1$.
The least-point construction applied to $M$ gives the left component
coordinates with bound $\beta$. The sum construction therefore gives
$z_q$ with bound $D_q$. Applying the block lemma to the outer natural
coordinate gives $r_{q+1}$ with bound $A_{q+1}$. Positivity is preserved
throughout by addition, successor, and multiplication by $\omega$.
This proves strict monotonicity and the stated bounds by induction.
For linear orders a strictly increasing map is injective. It reflects
strict order as well: if $r_q(x)<r_q(y)$ but $x\not<y$, then $y\le x$,
so monotonicity would give $r_q(y)\le r_q(x)$, a contradiction. Reflection
of non-strict order follows by taking the negation of the reversed strict
inequality. Hence the map is an order embedding.
\end{proof}

\begin{theorem}[Evaluation on the canonical invariant]
\label{hir:disjoint-invariant-ordinal-evaluation}
\uses{hir:disjoint-invariant-ordinal-embedding,hir:canonical-centres}
\lean{TranscendentalNakamaye.Pending.HironakaDisjointInvariantOrdinal.actualValueCode,TranscendentalNakamaye.Pending.HironakaDisjointInvariantOrdinal.actualValueCode_lt_bound,TranscendentalNakamaye.Pending.HironakaDisjointInvariantOrdinal.actualValueCode_lt_iff}
\leanok
In positive ambient dimension, evaluate $r_q$ directly on the canonical
invariant of an original marked ideal, original tower, and point. The
result lies below $A_q$, and comparison of two such coordinates is
equivalent to comparison of the original invariant values.
\end{theorem}
\begin{proof}
Each evaluated canonical invariant has the original type
$\operatorname{ValTy}_q$. Apply the preceding range bound and order
reflection to these values. No domain condition is needed for these
type-theoretic conclusions.
\end{proof}

\begin{remark}
The binary coordinate retains every original chronological label. These
coordinates quantify the already defined value order. They make no
assertion of geometric strict decrease or termination of a new driver.
\end{remark}
